\documentclass[../main.tex]{subfiles}
\begin{document}

\chapter{VLIW}
\lessoninfo{
  video={Lesson 11, videos 1--9},
  textbook={H\&P \cite{hennessy2017} §3.7, Appendix H},
  keywords={very long instruction word, VLIW, out-of-order superscalar, in-order superscalar, NOP, backward compatibility, irregular code, code bloat, full predication, branch hint, stop bit, Itanium, EPIC, digital signal processor},
}

Lesson 10 showed that the compiler can raise the number of instructions a
processor completes per cycle by rewriting the program: it shortens chains of
dependences, places independent instructions next to each other, and removes
branches and calls.  In a very long instruction word (VLIW) processor this
idea is taken to its end point: the compiler alone decides which operations
execute in parallel and encodes the decision in the instructions, so that the
hardware needs none of the dependence-tracking machinery of the out-of-order
processors of Lessons 7--9.  This lesson compares VLIW with superscalar
designs, weighs its advantages against its disadvantages, and shows where the
approach has failed and where it has succeeded.

\section{Superscalar and VLIW processors}
\label{sec:l11-superscalar-vs-vliw}

A superscalar\index{superscalar} processor (Lesson 8) executes up to $N$
instructions per cycle.\source{Lesson 11, videos 1--3; H\&P §3.7;
\cite{fisher1983}.}  Designs differ in how many instructions they examine to
find $N$ that can execute together, and in how much of that search they leave
to the compiler (\cref{fig:l11-windows,tab:l11-compare}).  An out-of-order
design finds the parallelism of any program by itself, at the price of complex
hardware for renaming, window scheduling and the reorder buffer (Lessons
6--8).  An in-order design stops at the first instruction that must wait for
another, so it needs the compiler to place independent instructions next to
each other (Lesson 10).  A VLIW processor executes, without any checking, one
instruction of up to $N$ operations that the compiler has made independent.

\begin{figure}[htbp]
  \centering
  % What each design examines to find N = 4 operations for one cycle:
% out-of-order superscalar (window >> N), in-order superscalar (next N),
% VLIW (one long instruction).  Shaded = executed in this cycle.
\begin{tikzpicture}[x=1cm,y=1cm, font=\footnotesize\sffamily, >={Stealth[length=1.6mm]},
  ins/.style={draw, minimum width=0.55cm, minimum height=0.45cm, inner sep=0pt, fill=white},
  run/.style={ins, fill=ln-accent!20},
  lab/.style={anchor=east, align=right, text width=2.6cm},
  win/.style={ln-accent, line width=0.6pt}]
  % --- row labels
  \node[lab] at (2.9,0)    {\iflangja{アウトオブオーダ\\スーパースカラ}{out-of-order\\superscalar}};
  \node[lab] at (2.9,-1.4) {\iflangja{インオーダ\\スーパースカラ}{in-order\\superscalar}};
  \node[lab] at (2.9,-3.4) {VLIW};
  % --- program order arrow
  \draw[->, ln-muted] (3.2,-4.15) -- (6.2,-4.15)
    node[right, font=\scriptsize\sffamily] {\iflangja{プログラム順}{program order}};
  % --- row 1: out-of-order, window of 12, four scattered picks
  \foreach \i in {0,...,11} {
    \node[ins] (o\i) at (3.475+0.55*\i,0) {};
  }
  \foreach \i in {0,1,4,8} { \node[run] at (o\i) {}; }
  \node[ln-muted] at (10.25,0) {\dots};
  \draw[win] (3.2,0.35) -- (3.2,0.45) -- (9.8,0.45) -- (9.8,0.35);
  % --- row 2: in-order, window of 4, fourth depends on second
  \foreach \i in {0,...,11} {
    \node[ins] (n\i) at (3.475+0.55*\i,-1.4) {};
  }
  \foreach \i in {0,1,2} { \node[run] at (n\i) {}; }
  \node[ln-muted] at (10.25,-1.4) {\dots};
  \draw[win] (3.2,-1.05) -- (3.2,-0.95) -- (5.4,-0.95) -- (5.4,-1.05);
  \draw[->, ln-caution] (n1.south) to[bend right=50]
    node[below, font=\scriptsize\sffamily] {\iflangja{依存}{dependence}} (n3.south);
  % --- row 3: VLIW, long instructions of four slots
  \foreach \j in {0,1,2} {
    \draw[fill=white] (3.2+2.2*\j,-3.625) rectangle (5.4+2.2*\j,-3.175);
    \foreach \k in {1,2,3} {
      \draw[ln-rule] (3.2+2.2*\j+0.55*\k,-3.625) -- (3.2+2.2*\j+0.55*\k,-3.175);
    }
  }
  \fill[ln-accent!20] (3.2,-3.625) rectangle (5.4,-3.175);
  \draw (3.2,-3.625) rectangle (5.4,-3.175);
  \foreach \k in {1,2,3} { \draw (3.2+0.55*\k,-3.625) -- (3.2+0.55*\k,-3.175); }
  \draw[line width=0.8pt] (5.4,-3.625) -- (5.4,-3.175) (7.6,-3.625) -- (7.6,-3.175);
  \node[ln-muted] at (10.25,-3.4) {\dots};
  \draw[win] (3.2,-3.05) -- (3.2,-2.95) -- (5.4,-2.95) -- (5.4,-3.05);
  % --- window labels
  \node[font=\scriptsize\sffamily, text=ln-accent, anchor=south] at (6.5,0.42)
    {\iflangja{調べる範囲：$N$ よりはるかに多い命令}{examined: $\gg N$ instructions}};
  \node[font=\scriptsize\sffamily, text=ln-accent, anchor=south west, inner xsep=0pt] at (3.2,-0.95)
    {\iflangja{調べる範囲：次の $N$ 命令}{examined: the next $N$ instructions}};
  \node[font=\scriptsize\sffamily, text=ln-accent, anchor=south west, inner xsep=0pt] at (3.2,-2.95)
    {\iflangja{調べる範囲：1 命令（$N$ 個の操作）}{examined: 1 instruction ($N$ operations)}};
\end{tikzpicture}

  \caption{What each design examines to fill one cycle with up to $N=4$
    operations.  Shaded boxes are executed in this cycle; the in-order design
    stops at the first instruction that depends on an earlier one in its
    group.}
  \label{fig:l11-windows}
\end{figure}

\begin{definition}[Very long instruction word]
  \label{def:l11-vliw}
  A \term{very long instruction word}[VLIW]%
  \index{VLIW|see{very long instruction word}} (VLIW) architecture is an
  instruction set architecture in which each instruction consists of up to $N$
  operations, each equivalent to one instruction of an ordinary ISA, that are
  executed in parallel in the same cycle.  Choosing the operations that form
  an instruction, and ensuring that they are independent, is the task of the
  compiler; the processor performs no dependence checking.
\end{definition}

An instruction that carries $N$ operations is about $N$ times as long as an
ordinary one, which gives the architecture its name.  Since it does the work
of up to $N$ ordinary instructions, a VLIW processor that completes one
instruction per cycle is comparable to a superscalar processor with an IPC of
up to~$N$ (Lesson 6).\caution{Instruction counts and CPI of a VLIW program
are not comparable with those of an ordinary ISA; compare operations per
cycle.}

\begin{table}[htbp]
  \centering\small
  \caption{Superscalar and VLIW processors that execute up to $N$ operations
    per cycle.}
  \label{tab:l11-compare}
  \begin{tabular}{@{}lccc@{}}
    \toprule
                       & Out-of-order          & In-order              & VLIW \\
    \midrule
    Executes per cycle & $\le N$ instructions  & $\le N$ instructions  & 1 instr.\ ($\le N$ ops) \\
    Examines           & $\gg N$ instructions  & next $N$ instructions & 1 instruction \\
    Hardware           & complex               & simpler               & simplest \\
    Compiler           & can help              & must help             & does all the work \\
    \bottomrule
  \end{tabular}
\end{table}

\begin{example}[Packing operations into instructions]
  \label{ex:l11-packing}
  The fragment below adds four array elements and stores the sum.  The
  comments list its RAW dependences (Lesson 3); the four loads only read R8
  and are independent of each other.
  \begin{lstlisting}[language={[x86masm]Assembler}, numbers=none]
I1: LW  R1, 0(R8)
I2: LW  R2, 4(R8)
I3: LW  R3, 8(R8)
I4: LW  R4, 12(R8)
I5: ADD R5, R1, R2    ; RAW with I1, I2
I6: ADD R6, R3, R4    ; RAW with I3, I4
I7: ADD R7, R5, R6    ; RAW with I5, I6
I8: SW  R7, 16(R8)    ; RAW with I7
\end{lstlisting}
  A VLIW compiler takes the operations in program order and places each one
  in the earliest instruction that follows all the operations it depends on
  and still has a free slot; slots left over are filled with NOPs
  (\cref{tab:l11-packing}).  With four slots the eight operations occupy
  four instructions, hence four cycles, where an ordinary ISA needs eight
  instructions; with two slots they occupy five.
\end{example}

\begin{table}[htbp]
  \centering\small
  \caption{The fragment of \cref{ex:l11-packing} compiled for VLIW ISAs with
    four and with two operation slots.  Each row is one instruction; sizes
    assume 32-bit slots.}
  \label{tab:l11-packing}
  \begin{tabular}{@{}cll@{}}
    \toprule
    Instruction & 4 slots            & 2 slots \\
    \midrule
    1           & I1, I2, I3, I4     & I1, I2 \\
    2           & I5, I6, NOP, NOP   & I3, I4 \\
    3           & I7, NOP, NOP, NOP  & I5, I6 \\
    4           & I8, NOP, NOP, NOP  & I7, NOP \\
    5           & ---                & I8, NOP \\
    \midrule
    NOP slots   & 8 of 16            & 2 of 10 \\
    Code size   & 512~bits           & 320~bits \\
    \bottomrule
  \end{tabular}
\end{table}

\section{Advantages and disadvantages}
\label{sec:l11-good-bad}

Leaving the search for parallelism to the compiler has three
advantages.\source{Lesson 11, videos 4--5; H\&P §3.7, Appendix H.}
\begin{itemize}
  \item \textbf{The compiler has plenty of time.}  An out-of-order processor
    must find independent instructions within a fraction of every clock cycle
    and only in its window; the compiler runs once and can analyze the whole
    program.
  \item \textbf{The hardware is simpler.}  Without renaming, window scheduling
    and a reorder buffer, a VLIW processor needs less area and power for the
    same operations per cycle, so it can be more energy efficient.
  \item \textbf{Loops and regular code work well.}  When arrays are processed
    in loops, the dependences are few and visible at compile time, and loop
    unrolling\index{loop unrolling} (Lesson 10) removes the branches between
    iterations, so that their independent operations fill the slots.
\end{itemize}
It also has four disadvantages.
\begin{itemize}
  \item \textbf{Latencies are not known at compile time.}  The compiler
    groups operations as if each completes within its cycle.  If one takes
    longer, e.g.\ a load with a cache miss (Lesson 12), all later operations
    wait for it, even independent ones, whereas an out-of-order processor
    keeps executing instructions that do not depend on it.
  \item \textbf{Irregular code is hard.}  The compiler sees the program, not
    its execution.  When it cannot prove two operations independent---a store
    and a load through pointers that may address the same location, or
    operations on either side of a data-dependent branch---it must place them
    in different instructions, whereas an out-of-order processor knows the
    actual addresses and predicts the branches at run time (Lessons 4 and~9).
  \item \textbf{Code bloat.}  Whenever fewer than $N$ independent operations
    are available, NOPs fill the slots and the program grows: the four-slot
    version in \cref{tab:l11-packing} is half NOPs and occupies 512~bits,
    against 256~bits for eight ordinary 32-bit instructions.  Compaction
    (\cref{sec:l11-instructions}) reduces this overhead.
  \item \textbf{A new processor needs recompiled programs.}  The grouping of
    operations and the number of slots are fixed at compile time, which
    undermines backward compatibility (below).
\end{itemize}

\begin{definition}[Backward compatibility]
  A processor offers \term{backward compatibility}[後方互換性] when it runs,
  unmodified, the programs compiled for earlier processors of its family.
\end{definition}

For superscalar processors backward compatibility comes with a performance
gain: a wider out-of-order processor finds more parallelism in an \emph{old}
binary without recompilation.  A new VLIW processor with more slots either
cannot run an old VLIW binary at all, because the instruction format has
changed, or runs it with no more operations per cycle than the processor it
was compiled for.

\begin{example}[An old binary on a wider VLIW processor]
  \label{ex:l11-compat}
  A four-slot processor that still accepts the two-slot format executes the
  two-slot version of \cref{ex:l11-packing} in five cycles, as the old
  processor does; only the recompiled four-slot version takes four.  The eight
  ordinary instructions, counted the same way, take five cycles on a two-wide
  and four on a four-wide out-of-order processor, without recompilation.
\end{example}

\begin{keypoint}
  VLIW trades hardware complexity for dependence on the compiler.  It pays off
  when the code is regular and compiled for the exact processor; it loses when
  parallelism or operation latencies become known only at run time, or when
  old binaries must run faster on new processors.
\end{keypoint}

\section{VLIW instructions}
\label{sec:l11-instructions}

A VLIW instruction is a sequence of $N$ operation slots
(\cref{fig:l11-format}).\source{Lesson 11, video 6; H\&P Appendix H.}  Each
slot is encoded like an instruction of an ordinary ISA, with its own opcode
from the full set of ordinary opcodes and its own operands.

\begin{figure}[htbp]
  \centering
  % A VLIW instruction with N = 4 operation slots, each encoded like an
% ordinary instruction, decoded and executed in parallel.
\begin{tikzpicture}[x=1cm,y=1cm, font=\footnotesize\sffamily, >={Stealth[length=1.6mm]},
  unit/.style={draw, fill=ln-box, minimum width=1.9cm, minimum height=0.5cm, inner sep=1pt}]
  % --- the instruction word: four contiguous slots of 2.45 cm
  \foreach \s in {0,...,3} {
    \pgfmathsetmacro{\x}{2.45*\s}
    \fill[ln-accent!15] (\x,0) rectangle (\x+0.95,0.55);
    \foreach \d in {0.95,1.45,1.95} { \draw[ln-rule] (\x+\d,0) -- (\x+\d,0.55); }
    \node[font=\scriptsize\sffamily] at (\x+0.475,0.275) {\iflangja{{\tiny 命令コード}}{opcode}};
    \node[font=\scriptsize\sffamily] at (\x+1.2,0.275) {rd};
    \node[font=\scriptsize\sffamily] at (\x+1.7,0.275) {rs};
    \node[font=\scriptsize\sffamily] at (\x+2.2,0.275) {rt};
    \pgfmathtruncatemacro{\n}{\s+1}
    \node[font=\scriptsize\sffamily, text=ln-muted, anchor=south] at (\x+1.225,0.55)
      {\iflangja{スロット \n}{slot \n}};
    % decode and execute
    \node[unit] (d\s) at (\x+1.225,-0.9) {\iflangja{デコード}{decode}};
    \node[unit] (e\s) at (\x+1.225,-1.8) {\iflangja{実行}{execute}};
    \draw[->] (\x+1.225,0) -- (d\s);
    \draw[->] (d\s) -- (e\s);
    \draw[->] (e\s) -- (\x+1.225,-2.6);
  }
  \foreach \s in {1,2,3} { \draw[line width=0.9pt] (2.45*\s,0) -- (2.45*\s,0.55); }
  \draw[line width=0.9pt] (0,0) rectangle (9.8,0.55);
  % --- architectural state
  \draw[fill=white] (0,-2.6) rectangle (9.8,-3.1);
  \node at (4.9,-2.85) {\iflangja{レジスタとメモリ}{registers and memory}};
  % --- marker for the whole word
  \draw[ln-accent, line width=0.6pt] (0,1.0) -- (0,1.1) -- (9.8,1.1) -- (9.8,1.0);
  \node[font=\scriptsize\sffamily, text=ln-accent, anchor=south] at (4.9,1.1)
    {\iflangja{1 命令 $=$ $N$ 個の独立な操作}{one instruction $=$ $N$ independent operations}};
\end{tikzpicture}

  \caption{A VLIW instruction with four operation slots.  Each slot is decoded
    and executed independently of the others; no hardware checks for
    dependences between slots.}
  \label{fig:l11-format}
\end{figure}

The operations of one instruction must be independent: none may use the
result of another, and no two may write the same location.  The hardware does
not check this, so correct grouping is the responsibility of the compiler.
Operations in \emph{consecutive} instructions may depend on each other, as I5
depends on I1 in \cref{tab:l11-packing}.\caution{Independence is required
only within an instruction, not between instructions.}

The compiler exposes the parallelism with instruction scheduling, as in
\cref{ex:l11-packing}, helped by the other techniques of Lesson 10: tree
height reduction\index{tree height reduction} shortens dependence chains so
that more operations can share an instruction, and loop unrolling and function
call inlining\index{function call inlining} remove the branches, calls and
returns that would separate operations.  A VLIW processor has no renaming
hardware, so the compiler itself removes false dependences (Lesson 3) by
writing results into otherwise unused registers, of which the instruction set
provides many.

Branches limit the grouping, since the operations after a branch execute only
on one of its paths.  A VLIW instruction set therefore provides good support
for predication (Lesson 5), ideally full predication\index{full predication},
so that the compiler can if-convert\index{if-conversion} branches and group
the operations of both paths.  A branch that remains can carry a
\term{branch hint}[分岐ヒント], with which the compiler tells the processor
which way it expects the branch to go.

Fixed slots waste space whenever fewer than $N$ independent operations are
available.  Compacted instructions store the operations one after another
without NOPs, and a \term{stop bit}[ストップビット] marks the last operation
of each group that executes together.  In the four-slot program of
\cref{tab:l11-packing}, I4, I6, I7 and I8 carry a stop bit; at 33~bits per
operation the code occupies 264~bits instead of 512.

\section{VLIW processors in practice}
\label{sec:l11-examples}

Two kinds of processors show where the VLIW approach fails and where it
succeeds.\source{Lesson 11, videos 7--9; H\&P §3.7, Appendix H.}

\subsection{Itanium}
\label{sec:l11-itanium}

The most prominent attempt to bring VLIW-style instructions to general-purpose
computing was \term{Itanium}[Itanium], the processor family Intel built for
the IA-64 instruction set; its designers called the approach EPIC, explicitly
parallel instruction computing \cite{schlansker2000}.\margin{IA-64 uses stops in
the same way to end groups of independent instructions
(\cref{sec:l11-instructions}).}  The instruction set, however, accumulated many
additional features whose hardware became very complex, removing the main
advantage of VLIW.  The compiler still found
little parallelism in irregular code, and Itanium did not displace
out-of-order processors.

\subsection{Digital signal processors}
\label{sec:l11-dsp}

A \term{digital signal processor}[DSP]%
\index{DSP|see{digital signal processor}} (DSP) is a processor specialized
for digital signals such as audio, video and radio.  This
market plays to the strengths of VLIW (\cref{sec:l11-good-bad}).\margin{H\&P
lists the TI C6x DSPs as an example of the VLIW approach.}
\begin{itemize}
  \item Loops that apply the same computation to many samples dominate the
    code: regular code in which the compiler finds plenty of parallelism.
  \item Programs are usually compiled for the particular DSP of the device they run
    in, so backward compatibility of binaries matters little.
  \item DSPs sit in power-limited devices, where energy-efficient hardware
    pays off.
\end{itemize}

\begin{keypoint}[Lesson summary]
  An out-of-order superscalar processor examines many more than $N$
  instructions and needs no help from the compiler; an in-order one examines
  $N$ and needs that help; a VLIW processor examines one
  instruction of up to $N$ independent operations and relies on the compiler
  entirely.  VLIW hardware is simple and energy efficient and performs well on
  regular loop code, but binaries must be recompiled for each new processor,
  irregular code yields little parallelism, a cache miss on one load delays
  all later operations, and NOPs bloat the code unless it is compacted
  with stop bits.  Itanium
  failed in general-purpose computing; DSPs are where VLIW fits.  Lesson 12 turns from instructions per cycle to memory.
\end{keypoint}

\end{document}
